\chapter{Time-Domain Response}
\label{chap:time-response}
\chapterepigraph{``To every thing there is a season, and a time to every purpose under the heaven.''}{Ecclesiastes 3:1, King James Version}

\section{What is ringdown?}

When a binary black hole merges, the observed gravitational wave passes in sequence through the growing-amplitude binary motion, a strongly nonlinear merger, and a decaying post-merger waveform (Fig.~\ref{fig:ringdown-regimes}). We call \term{ringdown} the time interval in which the post-merger remnant can be approximated as a single stationary black hole plus a small perturbation. This is not synonymous with ``the part of the waveform after its maximum amplitude.'' Immediately after the peak, nonlinear merger effects can remain, while at sufficiently late times a \term{power-law decay} generated by continuous components can dominate.

In a regime dominated by isolated quasinormal-mode poles, the dimensionless \term{gravitational-wave strain} $h(t)$ can be represented schematically as
\begin{align}
h_{\mathrm{RD}}(t)
&=
2\re\sum_{n\in C}
A_n\rme^{-\rmi\omega_n(t-t_0)},
\label{eq:ringdown-waveform-model}
\\
\omega_n
&=
\Omega_n-\frac{\rmi}{\tau_n},
\qquad
\Omega_n>0,
\quad
\tau_n>0.
\label{eq:ringdown-frequency-decay}
\end{align}
Here $A_n$ is a complex amplitude containing both the source and the observation operation, $\Omega_n$ is the angular oscillation frequency, $\tau_n$ is the damping time, and $t_0$ is a reference time chosen to parametrize phase and amplitude. Equation~\eqref{eq:ringdown-waveform-model} is not the definition of ringdown; it is a \term{pole approximation} valid in part of the ringdown regime.

\begin{figure}[tbp]
\centering
\resizebox{\linewidth}{!}{\begin{tikzpicture}[
  scale=1.1,
  every node/.style={font=\scriptsize},
  >=Latex
]
% Binary, merger, and post-merger black hole
\fill[black] (1.00,2.45) circle (.15);
\fill[black] (1.55,2.45) circle (.15);
\draw[gray,-{Latex[length=2mm]},line width=.45pt]
  (1.28,2.74)
  arc[start angle=90,end angle=360,x radius=.47,y radius=.28];
\node[gray] at (1.28,2.04) {binary};

\draw[gray,-{Latex[length=2mm]}]
  (2.00,2.45)--(2.72,2.45);

\fill[black] (3.16,2.45) circle (.23);
\fill[black] (3.46,2.45) circle (.23);
\draw[red,line width=.55pt]
  (3.00,2.45)
  .. controls (3.18,2.83) and (3.50,2.83) ..
  (3.66,2.45);
\node[red] at (3.31,2.04) {merger};

\draw[gray,-{Latex[length=2mm]}]
  (3.86,2.45)--(4.58,2.45);

\fill[black] (5.02,2.45) circle (.25);
\draw[blue,line width=.65pt]
  (4.64,2.45)
  arc[start angle=180,end angle=0,x radius=.38,y radius=.19];
\draw[blue!65,line width=.45pt]
  (4.56,2.45)
  arc[start angle=180,end angle=0,x radius=.46,y radius=.28];
\node[blue] at (5.02,2.04) {post-merger};

% Domains of validity
\fill[red!7]
  (3.45,-.92) rectangle (4.45,1.12);
\fill[blue!5]
  (4.45,-.92) rectangle (11.15,1.12);

\draw[red,line width=1.4pt]
  (3.45,1.31)--(4.45,1.31);
\node[red,above] at (3.95,1.31)
  {nonlinear general relativity};

\draw[blue,line width=1.4pt]
  (4.45,1.31)--(11.15,1.31);
\node[blue,above] at (7.80,1.31)
  {linear black-hole perturbation theory};
\node[blue,below,align=center] at (7.80,1.31)
  {$g_{\mu\nu}=g^{\mathrm{BH}}_{\mu\nu}+\delta g_{\mu\nu}$,
   $|\delta g|\ll|g^{\mathrm{BH}}|$};

% Waveform axes
\draw[gray,-{Latex[length=2mm]}]
  (.18,0)--(11.42,0)
  node[right] {$t$};
\draw[gray,-{Latex[length=2mm]}]
  (.34,-.92)--(.34,1.02)
  node[above] {$h(t)$};

\def\tpeak{3.48}
\def\tlin{4.45}
\def\ttail{8.10}

% Match amplitude and phase at the joining points.
\pgfmathsetmacro{\Apeak}{.055*exp(.60*\tpeak)}
\pgfmathsetmacro{\phipeak}{74*\tpeak+25*\tpeak*\tpeak}
\pgfmathsetmacro{\Alin}{\Apeak*exp(-.10*(\tlin-\tpeak))}
\pgfmathsetmacro{\philin}{\phipeak+395*(\tlin-\tpeak)}
\pgfmathsetmacro{\Atail}{
  \Alin
  *exp(-.50*(\ttail-\tlin))
  *sin(\philin+420*(\ttail-\tlin))
}

% Inspiral
\draw[
  black,line width=.75pt,smooth,
  domain=.45:\tpeak,samples=220
]
  plot
  (\x,{
    .055*exp(.60*\x)
    *sin(74*\x+25*\x*\x)
  });

% Nonlinear merger
\draw[
  black,line width=.95pt,smooth,
  domain=\tpeak:\tlin,samples=100
]
  plot
  (\x,{
    \Apeak
    *exp(-.10*(\x-\tpeak))
    *sin(\phipeak+395*(\x-\tpeak))
  });

% Exponential envelope of the QNM ringdown
\draw[
  red!60,dashed,line width=.45pt,smooth,
  domain=\tlin:\ttail,samples=80
]
  plot
  (\x,{
    \Alin*exp(-.50*(\x-\tlin))
  });
\draw[
  red!60,dashed,line width=.45pt,smooth,
  domain=\tlin:\ttail,samples=80
]
  plot
  (\x,{
    -\Alin*exp(-.50*(\x-\tlin))
  });

% QNM-dominated ringdown
\draw[
  red,line width=1.05pt,smooth,
  domain=\tlin:\ttail,samples=220
]
  plot
  (\x,{
    \Alin
    *exp(-.50*(\x-\tlin))
    *sin(\philin+420*(\x-\tlin))
  });

% Late-time power-law tail
\draw[
  blue,line width=1.05pt,smooth,
  domain=\ttail:11.15,samples=100
]
  plot
  (\x,{
    \Atail/pow(1+.62*(\x-\ttail),2.4)
  });

% Transition times
\draw[gray,densely dotted]
  (3.48,-1.02)--(3.48,1.05);
\draw[blue,densely dashed]
  (4.45,-1.02)--(4.45,1.12);
\draw[gray,densely dotted]
  (8.10,-1.02)--(8.10,1.05);

\node[below,gray] at (3.48,-1.00)
  {$t_{\mathrm{peak}}$};
\node[below,blue] at (4.45,-1.00)
  {$t_{\mathrm{lin}}$};
\node[below,gray] at (8.10,-1.00)
  {$t_{\mathrm{tail}}$};

% Ringdown and tail
\draw[red,line width=1.0pt]
  (4.58,-1.46)--(7.96,-1.46);
\node[below,red] at (6.24,-1.46)
  {QNM-dominated ringdown};

\draw[blue,line width=1.0pt]
  (8.24,-1.46)--(11.03,-1.46);
\node[below,blue,align=center] at (9.64,-1.46)
  {tail: $h\sim t^{-p}$, $p>0$\\
   branch cut / continuous component};
\end{tikzpicture}
}
\caption{Schematic waveform from binary coalescence to late times. $t_{\mathrm{peak}}$ lies near the maximum amplitude, $t_{\mathrm{lin}}$ is the time at which linear perturbation theory becomes accurate to a specified tolerance, and $t_{\mathrm{tail}}$ is the time at which the power-law decay from the continuous component can no longer be neglected. The latter two are not universal boundaries; they depend on the source, observable, and required accuracy.}
\label{fig:ringdown-regimes}
\end{figure}

It is therefore not generally correct either to fix the start of ringdown at $t_{\mathrm{peak}}$ or to assume that the waveform is exactly represented by a finite number of QNMs. Below we define this time-domain regime through the inverse transform of the causal Green function and ask when the \term{pole approximation} is justified.

\section{Moving the initial-value problem to frequency space}

Let the master-field equation be
\begin{equation}
\left(\partial_t^2+H\right)\Psi(t)
=
S(t).
\label{eq:time-response-wave}
\end{equation}
Denote the \term{initial data} at $t=0$ by $\Psi_0$ and $\Pi_0=\partial_t\Psi(0)$. Using the \term{one-sided Fourier--Laplace transform}
\begin{equation}
\widehat\Psi(\omega)
=
\int_0^\infty \rmd t\,
\rme^{+\rmi\omega t}\Psi(t),
\qquad
\im\omega>0,
\label{eq:one-sided-transform}
\end{equation}
integration by parts gives
\begin{equation}
\left(H-\omega^2\right)\widehat\Psi(\omega)
=
\widehat S(\omega)
+\Pi_0
-\rmi\omega\Psi_0.
\label{eq:initial-data-source}
\end{equation}
Thus the \term{initial data} enter the Green operator as an \term{effective source}:
\begin{equation}
\widehat\Psi(\omega)
=
\vG(\omega)
\left[
\widehat S(\omega)+\Pi_0-\rmi\omega\Psi_0
\right].
\label{eq:initial-response-resolvent}
\end{equation}
This single line already shows that QNM excitation depends on the \term{initial data}.

\section{The contour separates the waveform}

For $t>0$, the field is recovered from an upper-half-plane contour $C_+$ as
\begin{equation}
\Psi(t)
=
\int_{C_+}
\frac{\rmd\omega}{2\pi}\,
\rme^{-\rmi\omega t}
\widehat\Psi(\omega).
\label{eq:time-inverse-contour}
\end{equation}
Deforming the contour downward gives, schematically,
\begin{equation}
\Psi(t)
=
\Psi_{\mathrm{arc}}(t)
+\Psi_{\mathrm{QNM}}(t)
+\Psi_{\mathrm{cut}}(t).
\label{eq:prompt-ringdown-tail}
\end{equation}
The large-frequency arc gives the \term{prompt response} immediately behind the wavefront; isolated poles give QNM-dominated ringdown; and the \term{discontinuity} across a \term{branch cut} gives the \term{late-time tail}. These are not separate solutions but different pieces of one contour representation \cite{Leaver:1986gd}.

The sum over simple poles is
\begin{equation}
\Psi_{\mathrm{QNM}}(t)
=
-\rmi
\sum_{n\in C}
\rme^{-\rmi\omega_n t}
A_n
\left[
\widehat S(\omega_n)+\Pi_0-\rmi\omega_n\Psi_0
\right].
\label{eq:qnm-time-sum}
\end{equation}
Here $C$ denotes the set of poles crossed during the deformation. The overall sign depends on the contour orientation, but the pairing of pole exponentials and residues does not.

\section{The domain of a QNM expansion}

Equation~\eqref{eq:qnm-time-sum} is not a \term{complete expansion} valid at all times. Keeping only the QNM sum requires the large-arc contribution to vanish for the observation time and position, the \term{branch-cut} contribution to be negligible, and the pole sum to converge. In particular, summing QNMs before the wavefront arrives cannot by itself reproduce \term{causality}.

At very late times, on the other hand, a branch point near the origin can generate a \term{power-law decay} that dominates even the pole closest to the real axis in the lower half-plane. Ringdown is therefore generally an \term{intermediate-time asymptotic regime}.

\begin{warningbox}
The fact that a waveform can be fitted by finitely many damped sinusoids is not the same as saying that the Green function is represented by those poles alone. Vary the start and end times of the fit, the treatment of continuous components, and the \term{noise model}, and test the stability of pole locations and amplitudes separately.
\end{warningbox}

\section{Discontinuity across a branch cut}

Let $\vG_+$ and $\vG_-$ be the boundary values on the two sides of a branch cut, and define their \term{discontinuity} by
\begin{equation}
\operatorname{Disc}\vG(\omega)
=
\vG_+(\omega)-\vG_-(\omega).
\label{eq:green-discontinuity}
\end{equation}
Writing the branch cut as $C_{\mathrm{cut}}$, the continuous contribution is
\begin{equation}
\Psi_{\mathrm{cut}}(t)
=
\int_{C_{\mathrm{cut}}}
\frac{\rmd\omega}{2\pi}\,
\rme^{-\rmi\omega t}
\operatorname{Disc}\vG(\omega)
\ket{S_{\mathrm{eff}}(\omega)}.
\label{eq:cut-time-response}
\end{equation}
Here $\ket{S_{\mathrm{eff}}}$ collects the terms on the right-hand side of Eq.~\eqref{eq:initial-data-source}. The late-time power law is determined by the \term{small-frequency expansion} of the discontinuity near the origin.

In a spacetime bounded by two horizons, the potential decays exponentially at both ends. One must therefore not assume the same branch-cut structure as in an asymptotically flat spacetime. The boundary conditions must be used anew to determine which \term{nonanalyticities} are physical and which point sequences are artifacts of \term{numerical discretization}.

\section{From source to observer}

Let $A(t)$ be the detector signal and write its frequency-domain response as
\begin{equation}
\vA(\omega)
=
\bra{O}\vG(\omega)\ket{S}.
\label{eq:observable-frequency-response}
\end{equation}
Pole positions do not depend on the source or observer, but whether a given pole is visible depends on $\bra{O}A_n\ket{S}$. This matrix element may vanish by symmetry, nearby poles may cancel each other, or a continuous contribution may mask a pole. The ``\term{mode content}'' of ringdown is not a property of the background alone.

\begin{principlebox}
A decomposition of the time-domain waveform is a contour decomposition of the frequency-domain Green function. The boundaries between \term{prompt response}, ringdown, and the \term{late-time tail} are not absolute, but the analytic structures that generate them can be distinguished.
\end{principlebox}

\section{Conclusion of this chapter}

Initial data and external forcing enter the same Green operator, and poles, branch cuts, and the high-frequency arc combine into one causal time-domain response. This closes the classical-field response theory developed in Part I. In Part II we move poles and continuous components that are awkward to handle on the real axis onto spectra defined along complex coordinates.

\begin{researchproblem}[A uniform approximation for poles and tails]
For a fixed $\ell$ in Schwarzschild spacetime, numerically evaluate both the sum over QNM poles and the low-frequency branch-cut integral from the same \term{initial data}. Construct a \term{uniform approximation} that predicts the crossover time between ringdown and the tail using not only the smallest damping rate but also the residues and the support of the initial data. The diagnostic criterion is that a single error norm should control the discrepancy from direct time evolution after the wavefront arrives.
\end{researchproblem}
